\documentclass{standalone}
\usepackage{circuitikz}
\usetikzlibrary{calc}
\definecolor{circuitnotemark}{HTML}{FE00FE}
\makeatletter
\pgfdeclareshape{reg3}{
  \anchor{center}{\pgfpointorigin}
  \anchor{text}{\pgfpointorigin}
  \anchor{north}{\pgfpoint{0cm}{0.62cm}}
  \anchor{south}{\pgfpoint{0cm}{-0.62cm}}
  \anchor{east}{\pgfpoint{0.95cm}{0cm}}
  \anchor{west}{\pgfpoint{-0.95cm}{0cm}}
  \anchor{pin 1}{\pgfpoint{-1.35cm}{0cm}}
  \anchor{bpin 1}{\pgfpoint{-0.95cm}{0cm}}
  \anchor{pin 2}{\pgfpoint{0cm}{-1.02cm}}
  \anchor{bpin 2}{\pgfpoint{0cm}{-0.62cm}}
  \anchor{pin 3}{\pgfpoint{1.35cm}{0cm}}
  \anchor{bpin 3}{\pgfpoint{0.95cm}{0cm}}
  \backgroundpath{
    \pgfpathrectanglecorners{\pgfpoint{-0.95cm}{-0.62cm}}{\pgfpoint{0.95cm}{0.62cm}}
    \pgfpathmoveto{\pgfpoint{-1.35cm}{0cm}}\pgfpathlineto{\pgfpoint{-0.95cm}{0cm}}
    \pgfpathmoveto{\pgfpoint{0cm}{-1.02cm}}\pgfpathlineto{\pgfpoint{0cm}{-0.62cm}}
    \pgfpathmoveto{\pgfpoint{1.35cm}{0cm}}\pgfpathlineto{\pgfpoint{0.95cm}{0cm}}
  }
}
\makeatother
\begin{document}
\begin{circuitikz}[american, line width=0.8pt]
\ctikzset{bipoles/length=1.2cm}
\ctikzset{grounds/scale=1.36}
\foreach \x in {0,2,4,6,8,10,12} {\foreach \y in {0,-2,-4,-6,-8,-10,-12} {\fill[gray, opacity=0.35] (\x,\y) circle (1.1pt);}}
\node[gray, font=\scriptsize] at (0,0.84) {1};
\node[gray, font=\scriptsize] at (2,0.84) {2};
\node[gray, font=\scriptsize] at (4,0.84) {3};
\node[gray, font=\scriptsize] at (6,0.84) {4};
\node[gray, font=\scriptsize] at (8,0.84) {5};
\node[gray, font=\scriptsize] at (10,0.84) {6};
\node[gray, font=\scriptsize] at (12,0.84) {7};
\node[gray, font=\scriptsize, anchor=east] at (-0.84,0) {a};
\node[gray, font=\scriptsize, anchor=east] at (-0.84,-2) {b};
\node[gray, font=\scriptsize, anchor=east] at (-0.84,-4) {c};
\node[gray, font=\scriptsize, anchor=east] at (-0.84,-6) {d};
\node[gray, font=\scriptsize, anchor=east] at (-0.84,-8) {e};
\node[gray, font=\scriptsize, anchor=east] at (-0.84,-10) {f};
\node[gray, font=\scriptsize, anchor=east] at (-0.84,-12) {g};
\coordinate (a2) at (2,0);
\coordinate (c4) at (6,-4);
\coordinate (c7) at (12,-4);
\coordinate (e7) at (12,-8);
\coordinate (e4) at (6,-8);
\coordinate (g4) at (6,-12);
\node[vcc] at (a2) {VCC}; % line 3
\node[reg3, draw, font=\scriptsize] (part-U1) at (c4) {}; % line 4
\node[font=\scriptsize, anchor=south] at (part-U1.north) {$\mathrm{LM35}$}; % line 4
\node[anchor=south, yshift=9pt] at (part-U1.north) {$U_{1}$}; % line 4
\node[circuitnotemark, font=\tiny, anchor=west, xshift=2pt, inner sep=0] at (part-U1.bpin 1) {X}; % line 4
\node[circuitnotemark, font=\tiny, anchor=west, yshift=2pt, inner sep=0] at (part-U1.bpin 2) {X}; % line 4
\node[circuitnotemark, font=\tiny, anchor=east, xshift=-2pt, inner sep=0] at (part-U1.bpin 3) {X}; % line 4
\node[ground] at (c7) {}; % line 9
\draw (e7) node[ocirc]{} node[above left]{OUT}; % line 10
\draw (e4) to[C, l_=$C_{1}$, a^=$100\,\mathrm{n}\mathrm{F}$] (g4); % line 11
\node[ground] at (g4) {}; % line 12
\draw (a2) |- (part-U1.pin 1); % line 14
\draw (part-U1.pin 3) -| (c7); % line 15
\draw (part-U1.pin 2) -- (e4); % line 16
\draw (e4) -- (e7); % line 16
\node[circ] at (e4) {};
\node[anchor=south west, inner sep=2pt, yshift=4pt, circuitnotemark, font=\LARGE\bfseries] (circuittitle) at (current bounding box.north west) {X};
\path (circuittitle.south west) rectangle ++(15.972,0.853);
\end{circuitikz}
\end{document}
